% =====================================================================
% 6.5 Custo para um erro prescrito
% =====================================================================
\subsection{Custo para um erro prescrito}

O ajuste da Seção~\ref{sec:convergencia} dá, para cada método, uma lei \(E(h)\approx C h^{p}\). Com ela, o passo que produz um erro alvo é \(h=(E/C)^{1/p}\). O custo é o número de avaliações de \(f\) em \(T=5~\mathrm{ms}\): \(N_f=s\,T/h\), com \(s=1\), 2 ou 4 avaliações por passo. A Tabela~\ref{tab:custo} traz o resultado. Algumas metas pedem passos fora da faixa medida, de \(0{,}25\) a 8~\(\mu\)s. A meta de \(10^{-6}\)~A no Euler, por exemplo, pede \(h\approx 1{,}5~\mathrm{ns}\). Esses valores são extrapolações e devem ser lidos como ordem de grandeza.

\begin{table}[htbp]
\centering
\caption{Avaliações de \(f\) estimadas para atingir um erro máximo de corrente no intervalo de 5~ms.}
\label{tab:custo}
\small
\begin{tabular}{@{}rrrr@{}}
\toprule
Erro alvo (A) & Euler & Euler melhorado & RK4 \\
\midrule
\(10^{-3}\) & \(4{,}87\times 10^{3}\) & \(5{,}55\times 10^{2}\) & \(2{,}23\times 10^{2}\) \\
\(10^{-4}\) & \(4{,}27\times 10^{4}\) & \(1{,}76\times 10^{3}\) & \(3{,}96\times 10^{2}\) \\
\(10^{-6}\) & \(3{,}29\times 10^{6}\) & \(1{,}75\times 10^{4}\) & \(1{,}25\times 10^{3}\) \\
\(10^{-8}\) & \(2{,}54\times 10^{8}\) & \(1{,}75\times 10^{5}\) & \(3{,}96\times 10^{3}\) \\
\bottomrule
\end{tabular}
\fonte{Elaborado pelos autores (2026).}
\end{table}

Para um erro de \(10^{-6}\)~A, o RK4 precisa de cerca de \(1{,}3\times 10^{3}\) avaliações, o Euler melhorado, de \(1{,}8\times 10^{4}\), e o Euler, de \(3{,}3\times 10^{6}\). Quatro avaliações por passo parecem caras quando se olha um passo só. Para um erro fixo, porém, o RK4 é o método mais barato, porque pode usar um passo muito maior. \textcite[seção~10.6]{aquino2020} chega à mesma conclusão. O Euler só compete se a tolerância for grosseira e se o passo couber na região de estabilidade. No circuito com \(Q=50\), essa segunda condição já proíbe passos acima de 2~\(\mu\)s.

O custo mostra, portanto, que a ordem do método vale mais que o número de avaliações por passo.
